\section{The graph and its invariant subspaces}
\label{sec:graph}

Write $n=\prod_{i=1}^t p_i^{k_i}$, where the primes are distinct and
$k_i\geq1$. Denote by $G_n$ the graph whose vertices are the nonzero
proper ideals of $\mathbb Z_n$, with distinct ideals adjacent if their
intersection is nonzero. We use the combinatorial Laplacian $L=D-A$;
the graph is \emph{Laplacian integral} when every eigenvalue of $L$ is an
integer, including when the graph is disconnected.

The ideal $(\prod_i p_i^{\alpha_i})$ is encoded by
$0\leq\alpha_i\leq k_i$. We exclude $(0,\ldots,0)$, which encodes the
whole ring, and $(k_1,\ldots,k_t)$, which encodes the zero ideal.
Ideal intersection is generated by the least common multiple. Therefore
two distinct vertices $\alpha,\beta$ are adjacent exactly when
\[
\max(\alpha_i,\beta_i)<k_i\quad\hbox{for some }i.
\]
Equivalently, putting $S(\alpha)=\{i:\alpha_i<k_i\}$, adjacency is
$S(\alpha)\cap S(\beta)\ne\varnothing$.

For nonempty $S\subseteq[t]$, the support class $\mathcal C_S$ is a
clique of size $\prod_{i\in S}k_i$, except that the full support has
size $u=\prod_i k_i-1$. Edges between distinct support classes are complete
or absent according as their supports intersect or are disjoint. The full
support consists of universal vertices. Remove them and call the resulting
graph $H$; thus $G_n=K_u\vee H$, with the evident interpretation for $u=0$.

\begin{lemma}[Lifting across the universal vertices]
\label{lem:join}
If $L(H)v=\lambda v$ and $\lambda\ne0$, extension of $v$ by zero on
the $u$ universal vertices is an eigenvector of $L(G_n)$ with eigenvalue
$\lambda+u$.
\end{lemma}
\begin{proof}
Since $L(H)$ is symmetric with zero row sums, $\lambda\ne0$ implies
$\sum v=0$. At a vertex of $H$, its degree increases by $u$ and all
new neighbor values are zero. At a universal vertex the Laplacian action
is $-\sum v=0$. These are exactly the claimed eigenvector equations.
\end{proof}

For $(k_1,k_2,k_3)=(a,a,b)$, $H$ has six support classes of sizes
$a,a,b,a^2,ab,ab$ on $1,2,3,12,13,23$, respectively. The partition into
\[
X=\mathcal C_1\cup\mathcal C_2,\quad Y=\mathcal C_3,\quad
Z=\mathcal C_{12},\quad W=\mathcal C_{13}\cup\mathcal C_{23}
\]
is equitable, and on functions constant on these cells the Laplacian acts by
\begin{equation}
\label{eq:symmetric}
B=\begin{pmatrix}
a^2+ab&0&-a^2&-ab\\
0&2ab&0&-2ab\\
-2a&0&2a+2ab&-2ab\\
-a&-b&-a^2&a^2+a+b
\end{pmatrix}.
\end{equation}
For example, an $X$ vertex has $a^2$ neighbors in $Z$ and $ab$ in $W$,
while a $W$ vertex has $a,b,a^2$ neighbors in $X,Y,Z$.
Internal neighbors cancel for cell-constant functions. The cell sizes are
$2a,b,a^2,2ab$, and multiplying $B$ on the left by their diagonal matrix
gives a symmetric matrix. Thus $B$ is similar to a real symmetric matrix.

The functions taking values $s,-s$ on $\mathcal C_1,\mathcal C_2$,
$t,-t$ on $\mathcal C_{13},\mathcal C_{23}$, and zero elsewhere
form a second invariant subspace, with matrix
\begin{equation}
\label{eq:antisymmetric}
B_-=\begin{pmatrix}a^2+ab&-ab\\-a&a^2+a+b+2ab\end{pmatrix}.
\end{equation}
These functions have sum zero. Their eigenvalues therefore also lift by
$u=a^2b-1$, directly by the proof of Lemma~\ref{lem:join}.

Expanding~\eqref{eq:symmetric} gives $\det(xI-B)=x f_{a,b}(x)$, where
\begin{align}
f_{a,b}(x)&=x^3-E_1x^2+E_2x-E_3,\label{eq:cubic}\\
E_1&=2a^2+5ab+3a+b,\notag\\
E_2&=a^4+7a^3b+3a^3+8a^2b^2+11a^2b+2a^2+3ab^2+2ab,\notag\\
E_3&=2a^2b(a+b+1)(a^2+2ab+2a+b).\notag
\end{align}
Its roots are real and nonzero: reality follows from the similarity above,
and $E_3>0$. Every root $\lambda$ gives an eigenvalue $\lambda+u$ of
$G_n$. A noninteger root of $f_{a,b}$ therefore proves nonintegrality.
